% ============================================================
% Capítulo 4 — Dataset CUHK-PEDES
% ============================================================
\chapter{El dataset CUHK-PEDES}
\label{cap:dataset}

Este capítulo describe el dataset con el que se entrena y evalúa
el modelo: qué contiene, qué formato tiene en disco, cómo lo
carga el proyecto y qué decisiones de preprocesamiento se
aplican.

\section{El dataset en cifras}

\textbf{CUHK-PEDES} (\emph{Person Description in the Wild})
\cite{li2017personsearch} es el benchmark estándar para
\reid{} por texto. La versión \emph{original} tiene 40\,206
imágenes y 80\,440 descripciones (2 por imagen). Este proyecto
soporta dos variantes:

\begin{table}[H]
\centering
\caption{Diferencias entre la versión original y la versión
HuggingFace de CUHK-PEDES.}
\label{tab:cuhk-versions}
\begin{tabularx}{\textwidth}{lXX}
\toprule
 & \textbf{Original (paper)} & \textbf{HuggingFace} \\
\midrule
Total imágenes
  & 40\,206
  & 34\,042 \footnotesize(este proyecto) \\
Captions por imagen
  & 2
  & \textasciitilde 7 (1 por fila) \\
Identidades (pids)
  & 13\,003
  & derivadas del hash de imagen \\
Splits
  & train / val / test
  & derivados: 75\% / 12,5\% / 12,5\% \\
Resolución
  & variable (\textasciitilde 256×128+)
  & 128×128 thumbnails \\
Acceso
  & Email a Tong Xiao
  & Inmediato desde HuggingFace \\
Top-1 esperado
  & \textasciitilde 54\% (paper)
  & \textasciitilde 45-52\% \\
\bottomrule
\end{tabularx}
\end{table}

\noindent\textbf{Recomendación}: si se pueden obtener los datos
originales por email (gratis para investigación), se obtienen
mejores números. Si no, la versión HF funciona perfectamente
para validar el pipeline.

\section{Formato en disco}

Después de ejecutar \file{scripts/download_hf_dataset.py}, la
versión HF queda en \file{data/CUHK-PEDES-HF/} con esta
estructura:

\begin{verbatim}
data/CUHK-PEDES-HF/
|---- raw.parquet               # snapshot de HuggingFace (550 MB)
|---- reid_raw.json             # anotaciones en formato "estándar"
|__-- imgs/
    |-- imgs/all/
        |-- 7138159738431296927.jpg
        |-- 3825926175112857003.jpg
        `-- ... (34042 imágenes en total)
\end{verbatim}

\subsection{El archivo \texttt{reid\_raw.json}}

Es la \emph{fuente de verdad} del dataset. Cada elemento es un
diccionario con estos campos:

\begin{minted}[basicstyle=\ttfamily\small, frame=single]{python}
{
  "file_path": "all/7138159738431296927.jpg",
  "captions": [
    "A pedestrian with dark hair is wearing red and white shoes, "
    "a black hooded sweatshirt, and black pants.",
    "The person has short black hair and is wearing black pants, "
    "a long sleeve black top, and red sneakers.",
    ...
  ],
  "split": "train",
  "processed_tokens": [],
  "id": 0
}
\end{minted}

\noindent Los campos obligatorios son sólo \texttt{file\_path} y
\texttt{captions}. Los demás se completan en tiempo de carga.

\section{Cómo se cargan los datos}

El proyecto tiene dos clases de carga según la fuente:

\begin{itemize}
  \item \file{datasets/cuhkpedes.py}: \texttt{CUHKPEDES} para
        datos originales.
  \item \file{datasets/cuhkpedes_hf.py}: \texttt{CUHKPEDESHF}
        para datos de HuggingFace.
\end{itemize}

Ambas clases exponen la misma interfaz:

\begin{minted}[basicstyle=\ttfamily\small, frame=single]{python}
ds.train        # lista de (pid, image_id, img_path, caption)
ds.train_id_container    # set de pids en train
ds.val          # dict con image_pids, img_paths, caption_pids, captions
ds.test         # dict idem para test
\end{minted}

\subsection{El loader HF: \texttt{CUHKPEDESHF}}

Esta clase (\file{datasets/cuhkpedes_hf.py}) tiene tres pasos
clave:

\begin{enumerate}
  \item \textbf{Asignación de pids}: como la versión HF no
        incluye IDs de persona, los deriva del \emph{hash de la
        imagen}. Las imágenes idénticas (mismas captions)
        comparten el mismo pid.
  \item \textbf{División train/val/test}: por proporción de IDs:
        75\% / 12,5\% / 12,5\%. No usa el campo
        \texttt{split} del JSON.
  \item \textbf{Procesamiento de anotaciones}: igual que la
        versión original.
\end{enumerate}

\noindent El código clave está en \fline{datasets/cuhkpedes_hf.py}{62-71}:

\begin{minted}[numbers, firstnumber=62, linerange=62-71]{python}
def _build_pid_mapping(self):
    """Group entries by image_path (which is the identity key)."""
    self.path_to_id = {}
    next_id = 0
    for anno in self.annos:
        fp = anno['file_path']
        if fp not in self.path_to_id:
            self.path_to_id[fp] = next_id
            next_id += 1
    self.num_persons = next_id
\end{minted}

\subsection{El loader original: \texttt{CUHKPEDES}}

La clase \file{datasets/cuhkpedes.py} asume que las anotaciones
vienen con \texttt{split} explícito (train/val/test) y con
\texttt{id} entre 1 y 13\,003. Internamente resta 1 al id para
que empiece en 0 (más conveniente para PyTorch). El código
relevante está en \fline{datasets/cuhkpedes.py}{75-115}.

\section{El dataloader unificado}

\file{datasets/cuhkpedes_dataloader.py} es el punto de entrada
que decide qué loader usar y construye los \texttt{DataLoader}s
de PyTorch. Su función principal es
\texttt{build\_cuhkpedes\_dataloader(config)}.

\begin{minted}[numbers, firstnumber=22, linerange=22-86]{python}
def build_cuhkpedes_dataloader(config: dict = None):
    if config.dataset_source == 'huggingface':
        dataset_object = CUHKPEDESHF(config)
    else:
        dataset_object = CUHKPEDES(config)

    train_transform = T.Compose([
        T.ToTensor(),
        T.Normalize(mean=config.MHPV2_means, std=config.MHPV2_stds)
    ])
    inference_transform = T.Compose([
        T.ToTensor(),
        T.Normalize(mean=config.MHPV2_means, std=config.MHPV2_stds)
    ])
    train_num_classes = len(dataset_object.train_id_container)

    train_set = ImageTextDataset(
        dataset_object.train,
        train_transform,
        tokenizer_type=config.tokenizer_type,
        tokens_length_max=config.tokens_length_max
    )

    train_data_loader = DataLoader(
        train_set,
        batch_size=config.batch_size,
        shuffle=True,
        num_workers=config.num_workers,
        collate_fn=collate
    )
    ...
    return train_data_loader, inference_img_loader,
           inference_txt_loader, train_num_classes
\end{minted}

\noindent\textbf{Nota importante}: las transformaciones de
imagen \emph{sólo} convierten a tensor y normalizan. El resize a
$384 \times 128$ que el paper describe \textbf{no está}. Las
imágenes se quedan en su tamaño nativo ($512 \times 512$ después
del \texttt{place\_image\_on\_canvas}).

\section{La clase \texttt{ImageTextDataset}}

Es el \texttt{Dataset} de PyTorch que produce un batch mixto
imagen+texto. Vive en \file{datasets/bases.py}. Para cada índice:

\begin{minted}[numbers, firstnumber=124, linerange=124-151]{python}
def __getitem__(self, index):
    pid, image_id, img_path, caption = self.dataset[index]
    pid = torch.from_numpy(np.array([pid])).long()
    img = read_image(img_path)
    img = place_image_on_canvas(img)
    original_img = img.copy()
    if self.transform is not None:
        img = self.transform(img)
    # ... guarda también la imagen original sin normalizar
    tokens = self.tokenizer(caption)
    token_ids, orig_token_length = pad_tokens(tokens,
                                              self.tokens_length_max)

    ret = {'pids': pid,
           'image_ids': image_id,
           'img_paths': img_path,
           'preprocessed_images': img,
           'original_images': original_img,
           'token_ids': token_ids.to(torch.long),
           'orig_token_lengths': orig_token_length,
           'captions': caption}
    return ret
\end{minted}

\noindent Cada item del dataset devuelve un diccionario con 8
claves. El \texttt{collate} (\fline{utils/miscellaneous_utils.py}{189})
se encarga de apilar tensores y dejar las cadenas como listas.

\section{Preprocesamiento de imagen: \texttt{place\_image\_on\_canvas}}

Función definida en \file{datasets/bases.py}:

\begin{minted}[numbers, firstnumber=43, linerange=43-83]{python}
def place_image_on_canvas(image, new_size=(512, 512),
                          placement_strategy='random'):
    """Plaza una imagen en un canvas negro preservando aspect ratio."""
    if image.height > 512:
        aspect_ratio = image.width / image.height
        new_height, _ = new_size
        new_width = int(aspect_ratio * new_height)
        image = image.resize((new_width, new_height), Image.ANTIALIAS)

    canvas = Image.new('RGB', new_size, (0, 0, 0))

    if placement_strategy == 'center':
        x = (new_size[0] - image.width) // 2
        y = (new_size[1] - image.height) // 2
    elif placement_strategy == 'random':
        x = np.random.randint(0, max(1, new_size[0] - image.width))
        y = np.random.randint(0, max(1, new_size[1] - image.height))
    else:
        raise ValueError(...)

    canvas.paste(image, (x, y))
    return canvas
\end{minted}

\noindent\textbf{¿Por qué?} Las imágenes de CUHK-PEDES son
chicas ($\textasciitilde 128 \times 128$ en HF,
$\textasciitilde 256 \times 128$ en original). EfficientNet-B0
espera imágenes relativamente grandes ($224 \times 224$ o
más). Para no distorsionar el aspect ratio, se pega la imagen en
una posición aleatoria sobre un canvas negro de $512 \times 512$.
Esto introduce variación posicional como augmentation implícita.

\section{Preprocesamiento de texto}

\subsection*{Tokenización}

Hay tres tokenizers disponibles; el que se usa por defecto es
\textbf{BERT} (\texttt{bert-base-uncased}).

\begin{minted}[numbers, firstnumber=10, linerange=10-19]{python}
class BERTTokenizer(object):
    def __init__(self):
        super(BERTTokenizer, self).__init__()
        _ , tokenizer_class, pretrained_weights = (
            ppb.BertModel, ppb.BertTokenizer, 'bert-base-uncased')
        self.tokenizer = tokenizer_class.from_pretrained(
            pretrained_weights)

    def __call__(self, text: str = None) -> torch.LongTensor:
        tokens = self.tokenizer.encode(text)
        return tokens
\end{minted}

\noindent\textbf{Otros tokenizers} disponibles en
\file{datasets/simple_tokenizer.py} y
\file{datasets/tiktoken_tokenizer.py}, pero el config
(\fline{config.py}{41}) fija \texttt{tokenizer\_type='bert'} por
defecto.

\subsection*{Padding}

La función \texttt{pad\_tokens}
(\fline{utils/miscellaneous_utils.py}{86-107}) recorta o
padea a longitud fija:

\begin{minted}[numbers, firstnumber=86, linerange=86-107]{python}
def pad_tokens(tokens, tokens_length_max):
    tokens_length = len(tokens)
    tokens = torch.from_numpy(np.array(tokens)).view(-1)
                                  .contiguous().float()
    if tokens_length < tokens_length_max:
        zero_padding = torch.zeros(tokens_length_max - tokens_length)
        tokens = torch.cat([tokens, zero_padding], 0)
    else:
        tokens = tokens[:tokens_length_max]
        tokens_length = tokens_length_max
    return tokens, tokens_length
\end{minted}

\noindent\textbf{Tokens de padding = 0}. Esto es importante
porque el \texttt{nn.Embedding} del modelo tiene
\texttt{padding\_idx=0} (\fline{model/language_network.py}{22}),
lo que garantiza que el padding no afecte al embedding.

\section{Normalización de imagen}

El dataloader aplica \texttt{T.Normalize} con medias y desviaciones
de MHPV2 (no de ImageNet):

\begin{minted}[basicstyle=\ttfamily\small, frame=single]{python}
mean = [0.3555248975753784, 0.3295726776123047, 0.3115333914756775]
std  = [0.3482806384563446, 0.3339986503124237, 0.3296058773994446]
\end{minted}

\noindent definidos en \fline{config.py}{50-51}. \textbf{Esto es
un poco raro}: el paper usaría medias de ImageNet ($[0.485, 0.456,
0.406]$) que son las guardadas en \texttt{CUHKPEDES\_means}
(\fline{config.py}{52}). El código actual usa MHPV2 (que es
\emph{otro} dataset). Es un detalle a tener en cuenta al
reproducir resultados.

\section{El cuaderno de validación}

En \file{notebooks/01\_lab\_validate\_dataset.ipynb} hay un
notebook de Jupyter que verifica:

\begin{enumerate}
  \item Que el JSON existe y tiene los campos correctos.
  \item Que las imágenes referenciadas existen y son válidas.
  \item La distribución de longitud de captions.
  \item La distribución de splits.
  \item Que un forward pass del modelo funciona.
  \item Las pérdidas dan valores finitos.
  \item Una estimación de VRAM para elegir el \texttt{batch\_size}.
\end{enumerate}

\noindent Conviene correr este notebook \emph{antes} del primer
entrenamiento, especialmente si se usa una GPU diferente a la
del paper.

\section{Dataset augmentation}

Este proyecto \textbf{no aplica augmentation} explícita de
imagen más allá de:

\begin{itemize}
  \item Posición aleatoria en el canvas (\texttt{placement\_strategy='random'}).
  \item Normalización con medias MHPV2.
\end{itemize}

\noindent El paper original describe un pipeline más rico
(\fline{docs/paper.md}{245}):
\textit{redimensionamiento, recorte aleatorio y normalización}.
Este proyecto lo omite; podría añadirse en
\texttt{get\_transform} (\fline{utils/miscellaneous\_utils.py}{210}).

\section{Resumen del pipeline de datos}

\begin{figure}[H]
\centering
\begin{tikzpicture}[
    font=\sffamily\small,
    node distance=0.6cm,
    box/.style={
        rectangle, draw=black, thick, rounded corners=2pt,
        minimum height=0.9cm, minimum width=2.2cm, align=center,
        fill=blue!10
    },
    arr/.style={-{Stealth[length=2.5mm,width=2mm]}, thick}
]

\node[box] (json)   {reid\_raw.json};
\node[box, right=of json] (loader) {CUHKPEDES\\(HF)};
\node[box, right=of loader] (ds) {ImageTextDataset};
\node[box, right=of ds] (dl)  {DataLoader\\(collate)};
\node[box, right=of dl]  (out) {Batch\\\{preprocessed\_images,\\token\_ids,\\pids, ...\}};

\draw[arr] (json)   -- (loader);
\draw[arr] (loader) -- (ds);
\draw[arr] (ds)     -- (dl);
\draw[arr] (dl)     -- (out);

\node[below=0.2cm of loader, font=\scriptsize] {image\_pids, captions};
\node[below=0.2cm of ds,     font=\scriptsize] {place on canvas};
\node[below=0.2cm of dl,     font=\scriptsize] {shuffle, workers};
\end{tikzpicture}
\caption{Pipeline de datos desde el JSON hasta el batch que
consume \texttt{train.py}.}
\label{fig:pipeline-datos}
\end{figure}
